% GENERATED FILE. Edit canonical lessons, then run npm run generate:books.
% canonical-source-hash: fnv1a64:aaeed9946ce6ee06
% canonical-lessons: RU-C428-khostel, RU-C428-turbaza, RU-C428-nochleg, RU-C428-nochevat, RU-C428-bron

\chapter{A bed for the night}
\label{ch:ru-a-bed-for-the-night}
\hlchaptermodality{428}

\begin{chapteropening}
\textbf{By the end of this chapter:} \emph{I can find somewhere to sleep: a hostel, a holiday camp, a bed for the night, spending the night and a booking.}

Complete the last lesson of chapter 428: I can find somewhere to sleep: a hostel, a holiday camp, a bed for the night, spending the night and a booking.
\end{chapteropening}

\section[\texorpdfstring{\ru{хостел} (khóstel)}{хостел (khóstel)}]{\ru{хостел} (khóstel) — a hostel}

\label{lesson:RU-C428-khostel}



\begin{quote}
Before the new one: say the Russian for a captain, then the Russian for to sink.
\end{quote}

\subsection*{You'll want to know: \ru{хостел}}
\textbf{\ru{хостел}} (\emph{khóstel}) — \textquotedblleft{}a hostel\textquotedblright{}.

\begin{quote}
\ru{Студенты} \ru{жили} \ru{в} \ru{недорогом} \ru{хостеле}.
\end{quote}

\begin{quote}
\emph{Studénty zhíli v nedorogóm khóstele.}
\end{quote}

\begin{quote}
The students stayed in a cheap hostel.
\end{quote}

Masculine; the stress is on the first syllable.

\begin{cousinweb}
\emph{\ru{хостел}} is usually traced to English \emph{hostel}, which EtymDB traces through Old French to Latin \emph{hospitāle}, a place for guests. EtymDB does not list the Russian word.
\end{cousinweb}

\begin{grammarlens}[title={What you've built}]
One more word for telling, explaining and describing at B1.
\end{grammarlens}

\subsection*{Your turn}
\begin{itemize}
\raggedright
  \item \emph{Say it:} \emph{khóstel}
  \item \emph{Say it:} \emph{khóstel}, once more
  \item \emph{Say it:} a sentence of your own with \emph{khóstel}
  \item \emph{Recall:} say the Russian for a captain, then the Russian for to sink, then say \emph{khóstel} again
\end{itemize}

\begin{tcolorbox}[breakable,colback=teal!4,colframe=teal!35!black,title={Before you move on}]
What does \ru{хостел} mean? (\textquotedblleft{}A hostel\textquotedblright{}.) Say it once more.
\end{tcolorbox}
\section[\texorpdfstring{\ru{турбаза} (turbáza)}{турбаза (turbáza)}]{\ru{турбаза} (turbáza) — a holiday camp, a tourist centre}

\label{lesson:RU-C428-turbaza}



\begin{quote}
Before the new one: say the Russian for to sink, then the Russian for a hostel.
\end{quote}

\subsection*{You'll want to know: \ru{турбаза}}
\textbf{\ru{турбаза}} (\emph{turbáza}) — \textquotedblleft{}a holiday camp, a tourist centre\textquotedblright{}.

\begin{quote}
\ru{Летом} \ru{мы} \ru{отдыхали} \ru{на} \ru{турбазе} \ru{у} \ru{озера}.
\end{quote}

\begin{quote}
\emph{Létom my otdykháli na turbáze u ózera.}
\end{quote}

\begin{quote}
In summer we stayed at a holiday camp by the lake.
\end{quote}

Feminine; short for \emph{\ru{туристическая} \ru{база}}: simple cabins in the countryside.

\begin{cousinweb}
\emph{\ru{турбаза}} is a short form of \emph{\ru{туристическая} \ru{база}} 'a tourist base', from \emph{\ru{турист}} 'a tourist' and \emph{\ru{база}} 'a base'. EtymDB lists \emph{\ru{база}} and \emph{\ru{турист}} without a source; \emph{\ru{турист}} is usually traced through French \emph{touriste} and \emph{tour} to Latin \emph{tornāre} 'to turn on a lathe' (Lewis \& Short).
\end{cousinweb}

\begin{grammarlens}[title={What you've built}]
One more word for telling, explaining and describing at B1.
\end{grammarlens}

\subsection*{Your turn}
\begin{itemize}
\raggedright
  \item \emph{Say it:} \emph{turbáza}
  \item \emph{Say it:} \emph{turbáza}, once more
  \item \emph{Say it:} a sentence of your own with \emph{turbáza}
  \item \emph{Recall:} say the Russian for to sink, then the Russian for a hostel, then say \emph{turbáza} again
\end{itemize}

\begin{tcolorbox}[breakable,colback=teal!4,colframe=teal!35!black,title={Before you move on}]
What does \ru{турбаза} mean? (\textquotedblleft{}A holiday camp, a tourist centre\textquotedblright{}.) Say it once more.
\end{tcolorbox}
\section[\texorpdfstring{\ru{ночлег} (nochlég)}{ночлег (nochlég)}]{\ru{ночлег} (nochlég) — a bed for the night, a place to spend the night}

\label{lesson:RU-C428-nochleg}



\begin{quote}
Before the new one: say the Russian for a hostel, then the Russian for a holiday camp.
\end{quote}

\subsection*{You'll want to know: \ru{ночлег}}
\textbf{\ru{ночлег}} (\emph{nochlég}) — \textquotedblleft{}a bed for the night, a place to spend the night\textquotedblright{}.

\begin{quote}
\ru{Уже} \ru{темно}, \ru{а} \ru{ночлега} \ru{у} \ru{нас} \ru{ещё} \ru{нет}.
\end{quote}

\begin{quote}
\emph{Uzhé temnó, a nochléga u nas yeshchyó net.}
\end{quote}

\begin{quote}
It is already dark, and we still have nowhere to sleep.
\end{quote}

Masculine. \emph{\ru{Искать} \ru{ночлег}} — to look for a place to stay the night.

\begin{cousinweb}
\emph{\ru{ночлег}} is built from \emph{\ru{ночь}} 'night' and \emph{-\ru{лег}}, the root of \emph{\ru{лечь}} 'to lie down': a place to lie down for the night. EtymDB lists the word without a deeper source.
\end{cousinweb}

\begin{grammarlens}[title={What you've built}]
One more word for telling, explaining and describing at B1.
\end{grammarlens}

\subsection*{Your turn}
\begin{itemize}
\raggedright
  \item \emph{Say it:} \emph{nochlég}
  \item \emph{Say it:} \emph{nochlég}, once more
  \item \emph{Say it:} a sentence of your own with \emph{nochlég}
  \item \emph{Recall:} say the Russian for a hostel, then the Russian for a holiday camp, then say \emph{nochlég} again
\end{itemize}

\begin{tcolorbox}[breakable,colback=teal!4,colframe=teal!35!black,title={Before you move on}]
What does \ru{ночлег} mean? (\textquotedblleft{}A bed for the night, a place to spend the night\textquotedblright{}.) Say it once more.
\end{tcolorbox}
\section[\texorpdfstring{\ru{ночевать} (nochevát)}{ночевать (nochevát)}]{\ru{ночевать} (nochevát) — to spend the night}

\label{lesson:RU-C428-nochevat}



\begin{quote}
Before the new one: say the Russian for a holiday camp, then the Russian for a bed for the night.
\end{quote}

\subsection*{You'll want to know: \ru{ночевать}}
\textbf{\ru{ночевать}} (\emph{nochevát}) — \textquotedblleft{}to spend the night\textquotedblright{}.

\begin{quote}
\ru{Мы} \ru{ночевали} \ru{в} \ru{маленькой} \ru{гостинице}.
\end{quote}

\begin{quote}
\emph{My nocheváli v málenkoy gostínitse.}
\end{quote}

\begin{quote}
We spent the night in a small hotel.
\end{quote}

Imperfective; perfective \emph{\ru{переночевать}}. Present \emph{\ru{ночую}, \ru{ночуешь}}.

\begin{cousinweb}
\emph{\ru{ночевать}} is built on \emph{\ru{ночь}} 'night' with the verb ending \emph{-\ru{евать}}. EtymDB lists \emph{\ru{ночевать}} without a deeper source.
\end{cousinweb}

\begin{grammarlens}[title={What you've built}]
One more word for telling, explaining and describing at B1.
\end{grammarlens}

\subsection*{Your turn}
\begin{itemize}
\raggedright
  \item \emph{Say it:} \emph{nochevát}
  \item \emph{Say it:} \emph{nochevát}, once more
  \item \emph{Say it:} a sentence of your own with \emph{nochevát}
  \item \emph{Recall:} say the Russian for a holiday camp, then the Russian for a bed for the night, then say \emph{nochevát} again
\end{itemize}

\begin{tcolorbox}[breakable,colback=teal!4,colframe=teal!35!black,title={Before you move on}]
What does \ru{ночевать} mean? (\textquotedblleft{}To spend the night\textquotedblright{}.) Say it once more.
\end{tcolorbox}
\section[\texorpdfstring{\ru{бронь} (bron)}{бронь (bron)}]{\ru{бронь} (bron) — a booking, a reservation}

\label{lesson:RU-C428-bron}



\begin{quote}
Before the new one: say the Russian for a bed for the night, then the Russian for to spend the night.
\end{quote}

\subsection*{You'll want to know: \ru{бронь}}
\textbf{\ru{бронь}} (\emph{bron}) — \textquotedblleft{}a booking, a reservation\textquotedblright{}.

\begin{quote}
\ru{У} \ru{нас} \ru{есть} \ru{бронь} \ru{на} \ru{два} \ru{номера}.
\end{quote}

\begin{quote}
\emph{U nas yest bron na dva nómera.}
\end{quote}

\begin{quote}
We have a booking for two rooms.
\end{quote}

Feminine; colloquial for \emph{\ru{бронирование}}. Its cousin \emph{\ru{броня}} means armour.

\begin{cousinweb}
\emph{\ru{бронь}} goes with \emph{\ru{броня}} 'armour', which EtymDB traces through Old East Slavic \emph{brŭnja} to Proto-Slavic \emph{brŭnja} and on to Gothic \emph{brunjō} 'a breastplate': a booking shields a place for you.
\end{cousinweb}

\begin{grammarlens}[title={What you've built}]
One more word for telling, explaining and describing at B1.
\end{grammarlens}

\subsection*{Your turn}
\begin{itemize}
\raggedright
  \item \emph{Say it:} \emph{bron}
  \item \emph{Say it:} \emph{bron}, once more
  \item \emph{Say it:} a sentence of your own with \emph{bron}
  \item \emph{Recall:} say the Russian for a bed for the night, then the Russian for to spend the night, then say \emph{bron} again
\end{itemize}

\begin{tcolorbox}[breakable,colback=teal!4,colframe=teal!35!black,title={Before you move on}]
What does \ru{бронь} mean? (\textquotedblleft{}A booking, a reservation\textquotedblright{}.) Say it once more.
\end{tcolorbox}
